\documentclass[11pt]{article}
\usepackage[T1]{fontenc}
\usepackage{lmodern,amsmath,amssymb,amsthm}
\usepackage[margin=1in]{geometry}
\usepackage[hidelinks]{hyperref}
\newtheorem{theorem}{Theorem}
\newtheorem{lemma}[theorem]{Lemma}
\newtheorem{proposition}[theorem]{Proposition}
\newcommand{\len}{\operatorname{length}}
\newcommand{\Fitt}{\operatorname{Fitt}}
\newcommand{\im}{\operatorname{im}}
\newcommand{\OO}{\mathcal O}
\setlength{\emergencystretch}{1em}
\title{Finite coefficient lifts and integral jets of cyclotomic Gauss rows}
\author{\small\href{mailto:admin@cosmosgeek.org}{admin@cosmosgeek.org}\qquad
\href{mailto:b321014@iiit-bh.ac.in}{b321014@iiit-bh.ac.in}}
\date{Working paper for specialist review}
\begin{document}
\maketitle
\begin{abstract}
We bound the precision of a finite coefficient lift by the length of
an adjacent tensor product, the exponent of a cyclic critical quotient,
and the valuations of its tame Frobenius defects. Only one adjacent
module needs to have length strictly smaller than the critical precision.
We give equality examples, a two-branch order showing why an exponent
bound cannot replace a length bound, and an example in which a perfect
source pairing of fixed discriminant coexists with arbitrarily long
central relation loss. We compute every first-order relation on a cyclotomic Gauss row at level $p$ and prove that a compatible row admits unit-speed jets at every fixed precision at sufficiently deep levels. Joining two such rows can kill every speed. The arithmetic application requires an actual
finite lift of the original coefficient module and its class.
\end{abstract}

\section{A precision bound}
Let $\OO$ be a complete discrete valuation ring, $\pi$ a uniformizer,
$B=\OO/(\pi^L)$, and $D=\OO/(\pi^N)$ with $L\ge N>0$. All modules
are finitely generated and lengths are taken over $\OO$. A subscript
$D$ denotes tensoring over $B$ with $D$.
Let $C,V_1,V_2$ be $B$-modules with a multiplication map
\[
 V_1\otimes_B V_2\longrightarrow C.
\]
Write $J$ for its image, $M=C/J$, and $U=M_D$.
In particular $U=C_D/\im((V_1)_D\otimes_D(V_2)_D\to C_D)$
by right exactness. Lifting an image lattice or replacing a module
by a torsion-free quotient can change this equality.

\begin{theorem}\label{thm:precision}
Suppose $\Fitt_B^0(C)=0$ and $\len((V_i)_D)<N$ for at least one
$i\in\{1,2\}$. Put
$\kappa=\len((V_1)_D\otimes_D(V_2)_D)$.
Suppose there are $m_1,\ldots,m_s\in M$ with
$\pi^{a_i}m_i=0$. Write $T=\sum_i Bm_i$ and let $\bar T$ be
its image in $U$. If $U/\bar T$ is cyclic and killed by $\pi^b$
with $0\le b<N$, then
\[
 L\le\kappa+b+\sum_i a_i.
\]
\end{theorem}
\begin{proof}
For any $B$-module $V$, if $\len(V_D)<N$, its elementary divisors
all have length smaller than $N$. Indeed one summand of length at
least $N$ would contribute $N$ to the reduced length. Hence $V=V_D$
via reduction. Apply this to $M/T$: its reduction is $U/\bar T$,
of length at most $b$. Thus $\len(M)\le b+\sum a_i$. Also one
of $V_1,V_2$ equals its reduction, so their tensor product equals
the tensor of their two reductions. Therefore $\len(J)\le\kappa$.
Diagonal presentation of $C$ gives
$\Fitt_B^0(C)=(\pi^{\len(C)})B$, including when $C$ has a free
$B$-summand. Consequently $L\le\len C=\len J+\len M$.
\end{proof}

The strict cut $b<N$ cannot in general be dispensed with: for $T=0$
a single summand of arbitrary length $L\ge N$ reduces to $D$. The
total-length bound itself can be attained. Take $\kappa=0$ and
$C=M=\OO/(\pi^b)\oplus\bigoplus_i\OO/(\pi^{a_i})$,
with $m_i$ the corresponding basis vectors and $L=b+\sum a_i\ge N$.

\subsection{Tame valuations and scalar cocycles}
Let $K$ be a number field with residue characteristic $p$ for $\OO$,
and assume the coefficient ring is a $\mathbb Z_p$-algebra. Let a
continuous cocycle take values in $M(\delta)$, with scalar action
$\delta=\epsilon\Psi^{-1}$, where $\epsilon$ is the cyclotomic
character. At a finite place $v\nmid p$ assume $\Psi$ unramified,
write $F_v$ for arithmetic Frobenius and $q_v$ for the residue
cardinality. The restricted cocycle on inertia is a homomorphism.
It is zero on wild inertia and on the prime-to-$p$ part of tame inertia.
For a generator $\tau_v$ of pro-$p$ tame inertia write $m_v=c(\tau_v)$.
The relation $F_v\tau_vF_v^{-1}=\tau_v^{q_v}$ on the tame quotient
then gives
\[
 \delta(F_v)m_v=q_vm_v,\qquad (\Psi(F_v)-1)m_v=0.
\]
If a chosen lift of $\Psi(F_v)-1$ has valuation $a_v<L$, then
$\pi^{a_v}m_v=0$. Since inertia acts trivially, coboundaries vanish
on inertia and the coordinates specialize under reduction. In an
application with $U=D u_p+\bar T$ and $x^2u_p\in\bar T$, take
$b=2v_\pi(x)<N$. This gives the explicit precision budget of
Theorem~\ref{thm:precision} for any actual lift with that scalar class
and those exact coefficient modules. It assumes the zeroth Fitting
identity on $B$. One way to supply this identity is by base change
of a zeroth Fitting identity on the original ring, using its original
central module. An identity known only after passage to $D$ gives a
lower bound of $N$ in place of $L$.

A related estimate distinguishes individual elementary divisors.
Let $n_1\ge\cdots\ge n_h$ be the elementary lengths of $M$.
Assume that the residual functionals $M\to k=\OO/(\pi)$ inject
as classes into a $k$-space in which the subspace unramified outside
any subset $J$ of the tame places has dimension at most $1+|J|$.
Here the tame restriction of a functional $\lambda$ is its value on
$m_v$. For every $d\ge0$ the coordinate forms on summands of length
$>d$ vanish on the $m_v$ killed by $\pi^{a_v}$ with $a_v\le d$:
their projections to such summands are divisible by $\pi$. Thus
\[
 \#\{i:n_i>d\}\le 1+\#\{v:a_v>d\}.
\]
In decreasing order this bounds the $j$th length by $a_{j-1}$ for
$j\ge2$, after padding the list of valuations with zeros. This
conclusion uses the asserted residual injection and ramification
dimension bound as separate inputs.

\section{Tracking tame coordinates at a graph}
Let $S=\OO[[Y]]$ and suppose the coefficient is an $S$-module $M$.
If a tame scalar character has
$\Psi(F_v)=(1+Y)^{b_v}$ with $b_v\in\mathbb Z_p^\times$,
the quotient $((1+Y)^{b_v}-1)/Y$ is a unit of $S$. Thus $m_v\in M[Y]$.
In particular its reduction lies in $\im(M[Y]\to M/YM)$, the image
of the $Y$-torsion of the original module.

If $Y$ acts as a scalar $\rho$ of valuation $a>0$ on a finite
$\OO$-module with elementary lengths $n_i$, then
\[
 \len\im(M[Y]\to M/YM)
  =\sum_i \max\{0,a-|n_i-a|\}.
\]
On a cyclic summand of length $n$, the annihilator of $\rho$ is
$\pi^{\max(0,n-a)}M$ and the quotient has length $\min(n,a)$.
Their overlap gives the stated summand. A free $\OO$-summand contributes
zero. In particular the image can be nonzero for a module supported
on a single scalar graph even when its free quotient has trivial image.

\subsection{A first-order cut with nonsaturated product image}
Let $D=\OO/\pi^N$ and $A=D[\eta]/(\eta^2)$.
For a finite $A$-module $X$, set $X_0=X/\eta X$ and
$\Delta X=X[\eta]/\eta X$, naturally a submodule of $X_0$.
The periodic free resolution of $D$ over $A$ identifies
$\Delta X=\operatorname{Tor}_1^A(D,X)$.

Let $0\to J\to C\to M\to0$ be exact over $A$, with
$\Fitt_D^0(C_0)=0$ and $\len J_0\le a$. Assume $T\subseteq\Delta M$
and $M_0/T$ cyclic, killed by $\pi^{2a}$. Such a $T$ is supplied
by actual tame coordinates when the scalar character has unit tangent
at every retained tame place. Writing $H=\im(\Delta C\to\Delta M)$
gives
\[
 N\le\len C_0\le \min(N,2a)+\len J_0+\len H-\len(\Delta M/T)
 \le 3a+\len\Delta C .
\]
Indeed the kernel of $J_0\to C_0$ identifies with $\Delta M/H$ by
the Tor exact sequence. Substitution in
$\len C_0=\len M_0+\len J_0-\len(\Delta M/H)$ proves the bound.

Equality can coexist with a wholly invisible product in the critical
fibre. Put $N=3a$, $r=\pi^a$, $C=(A/r^2)e_1\oplus(A/r)e_2$,
$v=r e_1+e_2$, $J=A\eta v$ and $M=C/J$. Then $\Delta C=0$ since
$C$ is induced from $D$, $J_0$ has length $a$, and its image in $C_0$
is zero. The kernel of multiplication by $\eta$ on the constant
part of $M$ is $Dv$, so $\Delta M=Dv$ has length $a$ and $M_0/\Delta M$
is cyclic of length $2a$. Here $\len C_0=3a$. This is an example of
the coefficient estimate, and does not require a Galois realization.

\subsection{Length and exponent with two local rows}
Suppose a finite cut has two local relations. Let
$0\le m_i\le n$, $D=\OO/\pi^n$, and write
\[
 R(m;b,w)=\left(\bigoplus_{i=1}^s\OO/\pi^{m_i}\right)/(D b+D w),
 \qquad T_b=\sum_i D b_i e_i \subset R(m;b,w).
\]
A concrete source of this presentation is a scalar global
degree-one complex with one free coordinate $u$ and tame coordinates
$t_i$ of respective annihilators $d_i$, and a local map with rows
$\sum_i b_i t_i$, $u$, $\sum_i w_i t_i$. At precision $n$ put
$m_i=\min(n,v_\pi(d_i))$. The strict local functor in a finite module
$M$ consists of $u=0$ and the maps $R(m;b,w)\to M$.
For a universal quotient $U$ of $R(m;b,w)$ the obstruction
$\mathrm{Hom}_D(U,D)\to D^s$, given by the coordinates of the
$b_i e_i$, has image dual to the image of $T_b$ in $U$.
This uses exact duality over the principal Artin ring $D$.
A given arithmetic class must supply both rows and the stated quotient.

Write $v(0)=+\infty$. Put $S_0=\sum_i m_i$ and let $L(m;b,w)$ be
the minimum of the following numbers, omitting empty families:
\[
 \begin{gathered}
 S_0,\quad S_0-m_i+v(b_i),\quad S_0-m_i+v(w_i),\\
 S_0-m_i-m_j+v(b_iw_j-b_jw_i)\quad (i<j).
 \end{gathered}
\]
These are precisely the valuations of the possibly nonzero maximal
minors of the presentation $[\operatorname{diag}(\pi^{m_i})\mid b\mid w]$.
Smith reduction therefore gives $\len R(m;b,w)=L(m;b,w)$. To kill
$T_b$ add the individual columns $b_i e_i$. The $b$ column is then
redundant and the diagonal exponents become $k_i=\min(m_i,v(b_i))$.
Hence
\[
 \len T_b=L(m;b,w)-\sum_i k_i+
       \max\bigl(\{0\}\cup\{k_i-v(w_i):1\le i\le s\}\bigr).
\]
There is an equally explicit annihilator. Let $L=L(m;b,w)$ and let
$L_{-i}$ use the same formula with coordinate $i$ deleted. Adjoining
only $b_i e_i$ gives maximal minors of minimum valuation
$\min(L,v(b_i)+L_{-i})$ by expansion along that column. The cyclic
submodule $D b_i e_i$ thus has length $(L-v(b_i)-L_{-i})_+$. The
least $h\ge0$ such that $\pi^h T_b=0$ is the maximum of these lengths
(or zero for the empty set). This gives separate bounds for length and
annihilator even when the local slopes $b_i$ are nonunits.

\section{Three coefficient corners}
Consider a three-block matrix order over a commutative local ring $A$;
let $P,Q,C$ denote its $(1,2),(2,3),(1,3)$ corners. Use the convention
that $PQ$ denotes the sum of all central products. Write $R,S'$ for
the $(3,2),(2,1)$ corners. The same argument applies to coefficient
ideals of matrix blocks after choosing matrix units. Suppose $A/I=D$
and the relevant returns satisfy $RQ,PS'\subseteq I$ on the outer
diagonal corners. Then $CRQ,PS'C\subseteq IC$. Put
\[
 U_P=P/(CR+IP),\quad U_Q=Q/(S'C+IQ),\quad
 U_C=C/(PQ+IC).
\]
There is a right exact sequence
\[
 U_P\otimes_D U_Q\longrightarrow C/IC\longrightarrow U_C\longrightarrow0.
\]
Indeed the multiplication map kills the displayed relations, and its
image in $C/IC$ is exactly $(PQ+IC)/IC$. If $\Fitt_D^0(C/IC)=0$,
$U_P,U_Q$ are cyclic of length at most $a$, and $\len(U_C)\le2a$,
then $N\le3a$. The analogous bound with the actual tensor length and
the actual central length holds without cyclicity.

\subsection{Equality and a failure with exponent in place of length}
The bound $3a$ is attained. Take $A=\OO$, $I=(\pi^{3a})$,
$P=Q=\pi^a\OO$, $C=\OO$, $R=S'=\pi^{2a}\OO$, and the $(3,1)$
corner $\pi^{3a}\OO$. Diagonal corners are $A$. Ordinary matrix
multiplication preserves these ideals. The respective lengths of
$U_P,U_Q,U_C$ are $a,a,2a$, and $C/IC=D$.

For a two-branch example take odd residue characteristic and put
$A_j=\{(u,v)\in\OO^2:u-v\in\pi^j\OO\}$,
$A=A_{2a}$, $t=(\pi^{2a},-\pi^{2a})$, $I=tA$ and $a>0$.
The midpoint $(u+v)/2$ induces $A/I\simeq\OO/(\pi^{4a})$:
on $A$ it is multiplicative modulo $\pi^{4a}$ and its kernel
is $tA$. For the upper corners choose
\[
 P=Q=\pi^{2a} A_a,\quad C=A_0;
\]
for the lower corners choose $R=S'=\pi^{3a}A_0$ and
$(3,1)=\pi^{4a} A_0$. These ideals are closed under the required
products and the returns are contained in $I$. We have
$CR=S'C=\pi^{3a}A_0$, $IC=\pi^{2a}A_0$, and $PQ\subseteq IC$.
Thus $U_P\simeq U_Q\simeq\OO/(\pi^a)$ and
$U_C=C/IC\simeq(\OO/(\pi^{2a}))^2$. The Fitting ideal is zero
over $D$, and $\pi^{2a}U_C=0$, but $N=4a$. Replacing the length
hypothesis by the exponent hypothesis in the preceding bound fails.

\section{An integral monodromy example on a polarized local surface}
\label{sec:local-polarized-surface}
The position of a monodromy line relative to a specified block
decomposition matters even on the Tate module of a principally
polarized abelian surface. The following construction realizes the
full loss of primitivity in its projection, at arbitrary depth.

\paragraph{Construction.}
Fix $p=37$ and $N\ge1$. By Dirichlet's theorem choose a prime
$\ell>3$ with $\ell\equiv2+p^N\bmod p^{N+1}$ and $\ell\ne p$.
The ordinary case of
\href{https://www.numdam.org/articles/10.24033/asens.1183/}{Waterhouse,
\emph{Abelian varieties over finite fields}, Theorem 4.1}
gives an elliptic curve over $\mathbb F_\ell$ of trace $3$:
$|3|<2\sqrt\ell$ and $(3,\ell)=1$. Lift a nonsingular Weierstrass
model to obtain a good reduction curve $E_0/\mathbb Q_\ell$.
The roots $z,z'$ of $X^2-3X+\ell$ belong to $\mathbb Z_p$ and
are congruent to $2,1$ respectively, by Hensel's lemma. Since
$(1-z)(1-z')=\ell-2$ and $zz'=\ell$,
\[
 v_p(z'-1)=v_p(z-\ell)=N.
\]
The two integral Frobenius projectors on $T_p E_0$ split it into
its eigenlines; the alternating pairing pairs these lines perfectly.

Set $h=p^N$, and let $E_T/\mathbb Q_\ell$ be the Tate curve with
parameter $\ell^{h^2}$. In coordinates $1,4$ for $E_0$ and
$2,3$ for $E_T$, the product $A_0=E_0\times E_T$ has an integral
basis of $T=T_pA_0$ in which arithmetic Frobenius and a generator
of the pro-$p$ tame inertia image act by
\[
 F_0=\operatorname{diag}(z,\ell,1,z'),\qquad \tau_0=1+h^2E_{23}.
\]
Here $E_{ij}$ are matrix units. These formulas follow from the
Tate uniformization, taking a Frobenius lift fixing compatible
$p$-power roots of $\ell$. Choose the dual basis signs and inertia
generator so the alternating form has
$J_{14}=1$, $J_{23}=-1$, $J_{32}=1$, $J_{41}=-1$ and other
entries zero. We fix a generator of $\mathbb Z_p(1)$ in writing
this matrix.

Let $U=1+h^{-1}(E_{12}+E_{34})$. Then $U^tJU=J$.
On $UT$ the two matrices become $F=U^{-1}F_0U$ and
$\tau=U^{-1}\tau_0U$, with
\[
 F_{12}=(z-\ell)/h,\qquad F_{34}=(1-z')/h,\qquad
 \tau-1=h^2E_{23}-hE_{13}+hE_{24}-E_{14}.
\]
The unlisted entries of $F$ agree with $F_0$.
Both displayed off-diagonal Frobenius entries are units. The matrices
are integral invertible, so $UT$ is stable; wild and prime-to-$p$
tame inertia act trivially.

\paragraph{Polarization descent.}
Write $M=hU$. We have $h^2T\subset MT\subset T$ as $M$ and
$h^2 M^{-1}$ are integral, and $M^tJM=h^2J$.
Thus $D=MT/h^2T$ is a Galois-stable isotropic subgroup of
$A_0[h^2]$ of size $h^4$, the square root of the order of $A_0[h^2]$.
The quotient $\phi:A_0\to B=A_0/D$ has a principal polarization
pulling back to $h^2$ times the product polarization, by maximal
isotropy and polarization descent
(\href{https://www.jmilne.org/math/xnotes/AVs.pdf}{Milne,
\emph{Abelian Varieties}, Proposition 16.8}).
Under the rational Tate-module isomorphism the lattice of $B$ is
$h^{-2}MT=h^{-1}UT$. Rescaling by $h$ gives exactly the lattice
$UT$ and the alternating form $J$. Semistability and toric rank are
unchanged under isogeny; $B$ is semistable of toric rank one.

\paragraph{Integral position and block multiplication.}
On $T_pB$, written in this basis, we have
\[
 \tau-1=v w,\qquad v=(-1,h,0,0)^t,\quad w=(0,0,h,1).
\]
Both $v$ and $w$ are primitive, so the integral monodromy image is
the saturated line $\mathbb Z_p v$. Its projection to coordinates
$2,3$ is $h\mathbb Z_p e_2$. For reference, in the unramified
quadratic extension of $\mathbb Q_\ell$ the two forward adjacent
entries of $F^2$ are $(z^2-\ell^2)/h$ and $(1-z'^2)/h$.
They remain units since $p>2$ and $z\equiv\ell\not\equiv0\bmod p$.

This gives a direct check in a specified ambient matrix order.
Partition the coordinates into blocks $1\,|\,2\,|\,1$. In the order
with integral diagonal and upper blocks and lower blocks divisible
by $h$, the ideal generated by products around nontrivial block
cycles is $(h)$: each such cycle has a lower block, and
$E_{14}\cdot hE_{41}=hE_{11}$. Projecting each of the adjacent
entries of $F^2$ to its indicated single matrix coordinate and
inserting the middle matrix unit $E_{23}$ gives
\[
    ((F^2)_{12}E_{12})\,E_{23}\,((F^2)_{34}E_{34})
      =(F^2)_{12}(F^2)_{34}E_{14}.
\]
Its upper-right coefficient is a unit even modulo the block-cycle
ideal. The middle coordinate projection of the saturated monodromy
line therefore has arbitrarily large index in its saturation, while
this product generates the entire outer coordinate modulo $p^N$.
The order here is an explicitly chosen ambient order for a local
Galois action. No identification of that order with a global Hecke
algebra is used.

\subsection{The exterior square and the information retained by orientation}
Use the wedge basis $(12,13,14,23,24,34)$ and the symmetric pairing $B$
given by wedging into $\bigwedge^4$ with the chosen orientation. For
a matrix $g$ in a $1\mid2\mid1$ frame write $g_{14}=c$, $g_{1,2:3}=P$,
$g_{2:3,4}=Q$ and $g_{2:3,2:3}=M$. The block with rows $12,13$ and
columns $24,34$ of $\bigwedge^2g$ is $QP-cM$ by the $2$-by-$2$ minors.
For the Lie algebra action of $N$ this block is $-N_{14}I_2$. Consequently
the framed exterior-square monodromy recovers the outer entry directly,
including its integral divisibility.

An actual six-dimensional orthogonal similitude representation and a
chosen residual four-dimensional lift determine a unique lift through
the exterior-square map over a complete local ring with $2$ invertible.
It suffices to check square-zero extensions. Its differential
$\mathfrak{gl}_4\to\mathfrak{gso}(B)$ is an isomorphism: each
off-diagonal entry $x_{ij}$ is visible in a wedge entry with any third
index, and diagonal entries are recovered from paired sums by division
by two. These formulas give a left inverse; the two free Lie algebras
both have rank $16$, so the map is invertible over a local ring (reduce
modulo the maximal ideal). Correct the error of an arbitrary matrix
lift using the inverse differential. The uniqueness of the correction
preserves multiplication under lifting a representation. Passage to
the separated complete limit proves the assertion. This is the
integral local lifting statement, with the residual representation
and the split similitude frame specified.

Consider what is lost on retaining only characteristic polynomials.
Write $\delta=\det g$, $\sigma=\bigwedge^2g$, let $a_i$ and $b_i$ be
their elementary characteristic coefficients, and set
$P(\sigma)=\operatorname{pf}(B(\sigma-\delta\sigma^{-1}))$.
The pairing makes the matrix in this Pfaffian skew. With the indicated
ordering, the identities are
\[
 \begin{gathered}
 P(\sigma)=\delta a_1^2-a_3^2,\qquad b_1=a_2,\qquad
 b_2=a_1a_3-\delta,\\
 b_3=\delta a_1^2+a_3^2-2\delta a_2 .
 \end{gathered}
\]
Indeed on diagonal eigenvalues $x_i$ the Pfaffian is
$(x_1x_2-x_3x_4)(x_1x_3-x_2x_4)(x_1x_4-x_2x_3)$.
Expanding proves the formulas on diagonal matrices. Conjugation by a
similitude $h$ of multiplier $\nu$ multiplies the Pfaffian by
$\nu^3/\det h$, hence it is invariant on the image of $\mathrm{GL}_4$.
Density of diagonalizable matrices over characteristic zero and
injectivity into the rational coordinate ring give the identities
universally over the integers. Where $a_1$ is residually a unit and
$2$ is invertible, the oriented invariant determines
$a_1^2=(b_3+2\delta b_1+P(\sigma))/(2\delta)$; a specified residual
square root determines the lift over complete local rings.

For an explicit discrepancy take a rank-two integral representation
$T$ with determinant $\epsilon$, and a finite cyclic character $\chi$
of odd $p$-power order, over a coefficient ring containing its values.
Put $\rho_0=\epsilon\oplus T\oplus1$ and
$\rho_\chi=\epsilon\chi^{-1}\oplus(T\otimes\chi)\oplus\chi^{-1}$.
For $x=\zeta-1$, $\zeta$ a generator of the image of $\chi$, the
exterior-square trace difference on $g$ is
$\epsilon(g)(\chi(g)-\chi(g)^{-1})^2$ and generates $(x^2)$ as $g$
varies. In contrast the four-dimensional trace difference is
\[
 (\chi(g)-1)(\mathrm{tr}\,T(g)-\chi(g)^{-1}(\epsilon(g)+1)).
\]
If the mod-$p$ image contains a scalar $aI$ with $a\ne\pm1$ and that
element can be prescribed together with a primitive value of $\chi$,
these differences generate $(x)$. The same is true of the oriented
Pfaffian differences, since they are
\[
 \epsilon(g)^2
 \big((\mathrm{tr}\,T(g))^2-(\epsilon(g)+1)^2\big)
 (\chi(g)^2-\chi(g)^{-2}).
\]
At the chosen scalar the middle factor reduces to $-(a^2-1)^2$.
For example the stated simultaneous prescription holds for an
elliptic $T$ with full mod-$p$ image disjoint from the cyclic
character extension. Even oriented regular invariants do not recover
the monodromy lattice from the preceding local-surface example:
$\mathrm{diag}(1,t,t^{-1},1)$ fixes $F_0$ and scales
$\tau_0-1$ by $t^2$. For every finite tuple in the group generated by
these matrices the scaled tuple extends polynomially to $t=0$.
Simultaneous-conjugacy invariants take equal values there and at
$t=1$, while the specialized unramified representation has zero
monodromy. The distinction concerns integral matrix data relative to
the fixed frame. For the characteristic-zero invariant-theory context
see M.~Weidner, \href{https://arxiv.org/abs/1809.03644}{arXiv:1809.03644}.

\section{A source pairing and a central relation kernel}
There is also no bound on the following central relation loss from
the depth of a separate perfect source pairing over the same order.
Let $\OO=\mathbb Z_p$ with $p\ge7$, $\Lambda=\OO[[Y]]$ and
$R=\Lambda[t]/(t(t-p))$. In the basis $1,t$ define
$B(x,y)=\ell(xy)$ for $\ell(1)=0$, $\ell(t)=1$. Its Gram matrix is
$\bigl(\begin{smallmatrix}0&1\\1&p\end{smallmatrix}\bigr)$, and has unit determinant.
The Casimir element (the sum of each basis vector tensored with its
dual) is $1\otimes(t-p)+t\otimes1$. Its product specializes to
$-p$ under $t\mapsto0$, of valuation one.

For any $n\ge2$ put $d=p^n$, $D=\OO/d$ and $J=(d,Y)R$. Form the
ordinary block order over $R$ with scalar diagonal entries
(or a full $2$-by-$2$ middle block) and scalar entry ideals
\[
 \begin{pmatrix}
 R & R & J\\
 dYR & R & YR\\
 d^2R & dR & R
 \end{pmatrix}.
\]
The $1$-to-$3$ product ideal is $YR$. Closure follows, for example,
from $dJ\subset R$, $dYJ\subset YR$, $d^2YR\subset dYR$ and
$d^2R\subset dR$; the reverse-path products lie in $R$.
Its original central product quotient $J/YR$ is $R/YR$:
multiplication by $d$ on $R/YR$ is injective and $J=dR+YR$.
Under $R\to D[\eta]/(\eta^2)$ with $Y\mapsto\eta$ and $t\mapsto0$,
the quotient specializes to $D$, killed by $\eta$. Thus its
$\eta$-torsion maps onto a critical quotient of length $n$.
All $n$ share the same source order and the same source pairing.
The depth-one source Casimir cannot alone bound the length of this
central image. This is a construction with explicit modules and ideals;
an arithmetic application requires extra information comparing the
source with the original relation module.

\subsection{The depth supplied by a scalar eigenvector}
Let $D=\OO/\pi^B$, $A=D[\eta]/\eta^2$, and $t=t_0+\eta t_1$ in
a finite free $A$-module. Choose a basis over $\OO$ in which
$t_0=\pi^h e_1$, with $h<B$, and write
$t_1=a e_1+w$ with $w$ in the span of the remaining basis vectors.
Put $k=\min(B,v_\pi(a))$, and
$m=\min(h,\operatorname{content}(w))$, where the content of zero is
$+\infty$. Then the largest $L$ such that
$\mathrm{Ann}_A(t)\subseteq\pi^L A$ is
\[
 L=\max\{0,B-h-\max(0,m-k)\}.
\]
Indeed $x_0+\eta x_1$ kills $t$ exactly when $\pi^h x_0=0$,
$x_0w=0$ and $\pi^h x_1=-a x_0$. The least possible valuation of
$x_0$ allowing the last equation to be solved is
$\max(B-m,h-k)$. The last equation then bounds the valuation of
$x_1$ below by the displayed $L$. If $k\ge m$, the annihilator
$\eta\pi^{B-h}$ attains the bound (when $h=0$, the annihilator is zero).
Otherwise $a=\pi^k u$ with $u$ a unit. For $B-h-m+k\ge0$ use
$x_0=\pi^{B-m}$, $x_1=-u\pi^{B-h-m+k}$; when this exponent is
negative use $x_0=\pi^{h-k}$, $x_1=-u$.

Suppose an integral algebra $\mathcal H$ has original image $R$
on a free lattice and $t$ is an eigenvector after specialization,
with scalar character $\psi:\mathcal H\to A$. Every element in
$\ker(\mathcal H\to R)$ specializes to an operator killing $t$,
so $\psi$ descends to $R\to (\OO/\pi^L)[\eta]/\eta^2$.
The depth loss includes both the parallel and transverse first-order
components. In the central-module problem this step supplies a scalar
map. The class comparison and product-quotient hypotheses of
Theorem~\ref{thm:precision} remain arithmetic inputs.

\section{Integral first-order deformations of a cyclotomic Gauss row}
We now compute a local input to a finite-precision construction. The
objects in this section are explicit finite group-ring lattices. A
map from their integral relations to a central coefficient module
as above is an additional arithmetic hypothesis.

Fix an odd prime $p$, let $M=p^h$ for $h\ge1$,
$O_h=\mathbb Z_p[\zeta_{p^{h+1}}]$, and $R_h=O_h[T]/(T^M-1)$.
Write $K_h$ for the fraction field. Specify $r_h\in K_h[C_M]$ by its
values at the characters of $C_M$. Its augmentation is $1$. For a
character $\chi$ of exact order $p^j$, identify $T$ with the class of
$1+p$ and extend $\chi$ trivially on Teichm\"uller units; set
\[
 r_h(\chi)=\frac{(-1)^{j+1}}{p^j(p+1)}
   \sum_{u\in(\mathbb Z/p^{j+1})^\times}\chi(u)^{-1}
                   \zeta_{p^{j+1}}^{,u}.
\]
Take the roots of unity compatibly across levels. Finite Fourier inversion
defines $r_h$ and shows it projects to $r_k$ under $T\mapsto T$ for $k\le h$.
The character is primitive modulo $p^{j+1}$. Orthogonality gives the
product of its Gauss sum with the sum for $\chi^{-1}$ equal to
$\chi(-1)p^{j+1}=p^{j+1}$. The automorphism sending each root of
unity to its inverse exchanges these sums by $u\mapsto-u$; it preserves
the unique $p$-adic valuation. Hence
\begin{equation}\label{eq:gauss-size}
 v_p(r_h(\chi))=(1-j)/2.
\end{equation}

Let $L_h=R_h+R_hr_h$, $I_h=\ker(\mathrm{aug}:L_h\to O_h)$ and
$W_h=I_h/(T-1)I_h$. The elements $v=[T-1]$, $w=[r_h-1]$ generate
$W_h$. A normalized first-order map at precision $D=O_h/(\lambda^s)$,
$\lambda=\zeta_{p^{h+1}}-1$, is an $R_h$-linear homomorphism from $L_h$
to $D[\eta]/(\eta^2)$ reducing to augmentation and sending $1$ to $1$.
Here $T$ acts on the target by $1+u\eta$. Such a map exists exactly when a
functional $W_h\to D$ sends $v$ to $u$. Indeed $L_h=O_h\cdot1\oplus I_h$
and the coefficient of $\eta$ on $I_h$ must vanish on $(T-1)I_h$.

\subsection{All relations at the first level}
Set $h=1$, $e=p(p-1)$, and
\[
 a=(p-1)(p-3)/2,\qquad b=(p-1)^2/2,\qquad A=e-a .
\]
Then $\mathrm{aug}(R_1:r_1)=(\lambda^b)$ and
$W_1\simeq O_1/(\lambda^a)\oplus O_1/(\lambda^A)$.
At every depth $s$, the possible $u$ form $(\lambda^{k_s})D$, where
\[
 k_s=\min\{(s-a)_+,p-1\}+(s-A)_+ .
\]

Here are details. Put $N=1+T+\cdots+T^{p-1}$, $z=\zeta_{p^2}$ and
$G(T)=\sum_u T^{-\log_{1+p}\langle u\rangle}z^u$, where $\langle u\rangle$
is the principal-unit part of $u$ modulo $p^2$ and the sum ranges over
units. Then $r_1=N/p+G/[p(p+1)]$. Replacing $r_1$ by
$\epsilon=(N+G)/p=((p+1)-N)r_1$ preserves both the lattice and the
colon ideal: the multiplier is an augmentation-one unit with inverse
$(1+N)/(p+1)$. Compute over the degree-$p$ period field $F$ fixed by
Teichm\"uller automorphisms of $K_1$; the result extends flatly to $O_1$.
Let $t=1+\sum_{u\in\mu_{p-1}}z^u$ and let $\sigma(z)=z^{(1+p)^{-1}}$.
Then $t$ is a uniformizer of $F$, $v_\lambda(t)=p-1$, and
$\sigma(t)=t+ct^2+O(t^3)$ in $O_F/p$, $c\in\mathbb F_p^\times$.
Indeed in expanding $(1+\lambda)^u$ the power sums over the
Teichm\"uller residues first contribute a unit at degree $p-1$. For
$\sigma(t)-t$ the first term modulo $p$ is
$-\lambda^p\sum_u u(1+\lambda)^u$ up to terms of degree at least $2p$.
The first nonzero power sum in this expression is at degree $p-2$.
Thus the leading valuations are respectively $p-1$ and $2p-2$.
In particular $O_F/p=\mathbb F_p[t]/(t^p)$.

Set $x=T-1$ and write $p\epsilon=\sum A_i x^i$, with representatives
of degree $<p$. Its coefficients in $T$ are $\sigma^m(t)$,
$0\le m<p$, and $A_0=p$. For $i\ge1$ we have in $O_F/p$
\[
 A_i=t^{p-i}(\gamma_i+O(t)),\qquad \gamma_i=-c^{p-1-i}/i!.
\]
In fact $(\sigma-1)^jt=j!c^jt^{j+1}+O(t^{j+2})$ for $j\le p-2$.
Expand $\sigma^m(t)$ binomially, noting that the sum over $m$ of
$\binom mi\binom mj$ vanishes in $\mathbb F_p$ if $i+j<p-1$ and equals
$-1/(i!j!)$ if $i+j=p-1$.

Put $k=(p-1)/2$ and $J=(O_F[C_p]:\epsilon)$. For $B=\sum B_ix^i\in J$,
$B(p\epsilon)$ vanishes in $\mathbb F_p[t,x]/(t^p,x^p)$.
If $w_0=\min_i(v_t(B_i)+i)<k$, examine the coefficient of
$t^{w_0+p-l}x^l$ for $l=w_0+1,\ldots,2w_0+1$. Writing $b_i$ for
the coefficients of $t^{w_0-i}$ in $B_i$, we obtain
$\sum_{i=0}^{w_0} b_i\gamma_{l-i}=0$. After multiplying rows/columns
by units and rows by $l!$, the columns are the falling factorials
$(l)_i$ at distinct residues. The Vandermonde determinant is a unit.
This contradiction gives $v_t(B_i)+i\ge k$, hence $v_t(\mathrm{aug}(B))\ge k$.
Equality is attained: set $B_0=t^k$, $B_i=t^{k-i}U_i$ for $1\le i\le k$,
and the rest zero. The equations of $x$-degree $l\le k$ hold automatically.
For $k+1\le l\le2k$ divide by $t^{k+p-l}$. Arbitrary power-series lifts of
$A_{l-i}/t^{p-l+i}$ give a square system in the $U_i$ with invertible
constant coefficient matrix: after unit scalings it is evaluation of
$(l)_i$, $1\le i\le k$, and $l$ is a unit. Solving over
$\mathbb F_p[[t]]$ and restoring the powers of $t$ solves the original
congruences. Ambiguities of lifts disappear modulo $t^p$. Thus
$\mathrm{aug}(J)=t^k O_F$.

Write $C=B\epsilon$ in the integral group ring, and compute with
degree-$<p$ representatives. Let $d(B)=(C-B)'(1)$, where the derivative
is with respect to $T$. Polynomial division yields
$B(p\epsilon)=pC+((1+x)^p-1)Q$. Differentiation at $x=0$ gives
\[
 d(B)=B(1)\epsilon'(1)-Q(0).
\]
Here $v_t(\epsilon'(1))=-1$, by the formula for $A_1$. Moreover
$Q(0)\equiv[x^p] B(p\epsilon)\pmod p$, since every nonleading
coefficient of $(1+x)^p-1$ is divisible by $p$. Each product
$B_iA_{p-i}$ contributing has valuation at least $k$. It follows that
$d(B)\in t^{k-1}O_F$ and this valuation is attained when $B(1)$
generates $t^kO_F$.

For any $h$ the relations on $v,w$ are generated by
$(p^h,0)$ and $(-d(B),\mathrm{aug}(B))$ for $Br_h=C\in R_h$,
where $d(B)=(C-B)'(1)$ using canonical representatives. To see
this, express $I_h=xR_h+R_h(r_h-1)$. A syzygy has
$B(r_h-1)\in xR_h$, equivalently $Br_h=C\in R_h$, with
$C(1)=B(1)$. Dividing $C-B$ by $x$ gives the derivative after taking
augmentation; the ambiguity is a multiple of the group norm, of
augmentation $p^h$. Reducing syzygies modulo $x$ gives the claimed
presentation by right exactness. The same presentation applies to
$\epsilon$ at $h=1$. The determinant of $x$ on the rational
augmentation-zero space is $\pm p^h$, so $\len W_h=v_\lambda(p^h)$.
At $h=1$ the smallest relation valuation is $(k-1)(p-1)=a$.
The other invariant is therefore $A$. Projection of the relation
lattice onto its second coordinate has image $(\lambda^b)$. A triangular
basis has rows $(\lambda^{e-b},0)$ and $(u_0\lambda^a,\lambda^b)$
with $u_0$ a unit, because $a<b,e-b$. Solvability of the two
relations in $D$ gives the stated formula for $k_s$.

\subsection{Arbitrary fixed precision at large depth}
At every fixed integer $J>0$ there is $h_0(p,J)$ such that for every
$h\ge h_0$ the entire row $L_h$ admits a normalized first-order map
with $u=1$ modulo $p^J$. We give a proof because knowing a single
character value or a derivative of a rational representative does not
control all integral syzygies.

First we need a fact about bounded power series. In a complete
nonarchimedean field containing $\mathbb C_p$, let $v(p)=1$ and
$f(X)=\sum_i a_i X^i\ne0$ with $\mu=\inf_i v(a_i)>-\infty$.
For $S_j=\{\zeta-1:\zeta\text{ of exact order }p^j\}$ and
$d_j=\varphi(p^j)$,
\[
 \lim_{j\to\infty}j^{-1}\min_{x\in S_j}v(f(x))=0 .
\]
Indeed on the closed disc of valuation radius $1/d_j$ the Gauss valuation
$\mu_j=\min_i(v(a_i)+i/d_j)$ tends to $\mu$. The largest minimizing
index $l_j$ satisfies $l_j/d_j\le\mu_j-\mu\to0$. One-variable
Weierstrass preparation in the Tate algebra, after scaling by a primitive
$\zeta-1$, gives $l_j$ roots $\alpha_i$ in the disc with multiplicity and
a constant-valuation unit; on the disc
$v(f(x))=\mu_j-l_j/d_j+\sum_i v(x-\alpha_i)$.
Work in a splitting field as needed. For every root remove a nearest
point $x_i$ of $S_j$. At every remaining point
$v(x-\alpha_i)\le v(x-x_i)$. The sum over $x\ne x_i$ of $v(x-x_i)$ is
$v(\Phi'_{p^j}(1+x_i))=j-1/(p-1)$ by the cyclotomic derivative formula.
All valuations in this sum are nonnegative. Averaging over the remaining
at least $d_j-l_j$ points gives an upper bound
$\mu_j-l_j/d_j+l_j(j-1/(p-1))/(d_j-l_j)$ for the minimum, while
$\mu_j$ bounds it below. Division by $j$ proves the assertion.

Use valuations normalized by $v(p)=1$. Write $\beta_h$ for the
valuation of a generator of $\mathrm{aug}(R_h:r_h)$ and
$\tau_h=h-\beta_h$. Since $W_h$ has total normalized length $h$ and
$W_h/O_h v$ is the augmentation quotient of valuation $\beta_h$, the
order ideal of $v$ has valuation $\tau_h$. Equivalently $\tau_h$ is
the least valuation of an integral scalar $c$ with $c e_h\in L_h$,
where $e_h=(1+\cdots+T^{p^h-1})/p^h$. Indeed $c v=0$ in $W_h$ means
$x(c-y)=0$ for some $y\in I_h$, or $c-y=c e_h$.

We claim that both $\tau_h$ and $\beta_h$ tend to infinity.
A bounded subsequence of $\tau_h$ would give integral rows $A_h,B_h$
and a fixed integer $C$ with $p^C e_h=A_h+B_h r_h$ on that sequence.
A bounded subsequence of $\beta_h$ would instead give integral
$A_h=B_h r_h$ with $\mathrm{aug}(B_h)=p^C$. The following common
argument rules out both cases.

Take bounded sequences in $\mathbb C_p$, modulo sequences whose norm
tends to zero along a nonprincipal ultrafilter on the chosen levels.
The quotient is a valued field with norm the ultralimit of the norms;
a nonzero class has a bounded inverse on a set in the ultrafilter.
Let $K$ be its completion and embed $\mathbb C_p$ by constant sequences.
Project the integral $A_h,B_h$ to each fixed $C_{p^j}$. Coefficientwise
they are still bounded by one, hence they yield compatible integral
rows $A_j^*,B_j^*$ over $K$. Compatibility holds before passing
to the quotient for all sufficiently large $h$. The coefficients of
the $B_j^*$ on residues modulo $p^j$ define a bounded measure on
$\mathbb Z_p$: the projections give additivity, and integration extends
from locally constant functions by uniform approximation. Its binomial
moments define a bounded series
$f(X)=\sum_i(\int\binom zi\,d\mu(z))X^i$.
For every $p$-power root $\zeta$, uniform convergence of the binomial
expansion yields $f(\zeta-1)=B_j^*(\zeta)$ at the corresponding level.
By \eqref{eq:gauss-size} and integrality of $A_j^*$ at these characters,
$v(f(\zeta-1))\ge(j-1)/2$ on all $S_j$. The series lemma forces
$f=0$. Finite Fourier inversion makes $B_j^*=0$ at every level.
In the first case $A_j^*=p^C e_j$ contradicts integrality for $j>C$;
in the second case $\mathrm{aug}(B_j^*)=p^C$ is a contradiction.

Finally let $\alpha_h$ be the minimum normalized valuation of the
entries of all relations on $v,w$. For $h\ge k$, projection to level
$k$ preserves augmentation and maps the colon ideal into the lower
colon ideal after flat extension of scalars. The derivative of
a canonical group polynomial differs from the derivative of its
projection by a multiple of $p^k$, by grouping exponents modulo $p^k$.
At level $k$ choose one colon element whose augmentation generates that
ideal. Any projected upper relation subtracts an integral multiple of
this relation to eliminate its second coordinate; the multiplier has
valuation at least $\beta_h-\beta_k$. A relation with second coordinate
zero annihilates $v$, so its first coordinate has valuation at least
$\tau_k$. Since $\alpha_k\le\beta_k$ and $\tau_k\le k$, this proves
\[
 \alpha_h\ge\min(\beta_h-\beta_k+\alpha_k,\tau_k).
\]
Given $J$, first take $k$ with $\tau_k>J$ and then $h$ with
$\beta_h-\beta_k+\alpha_k>J$. Every entry is divisible by $p^J$.
Sending $v$ to $1$ and $w$ to $0$ now gives the desired functional.
No effective bound for $h_0$ follows from this compactness argument.

There is also an explicit limitation when the target precision varies
with $h$. Write $F=K_h^{\mu_{p-1}}$, $M=p^h$ and normalize
$v_F(p)=M$. Fourier inversion, or the defining sums, shows the row
coefficients lie in $F$. The automorphism $\sigma$ induced by
$\zeta\mapsto\zeta^{1+p}$ on those coefficients sends $r_h$ to $T r_h$.
Indeed on each primitive character stratum the change of additive
exponent multiplies the Gauss row by $T$, and the projectors onto
the strata are rational. Flat extension permits computing the two
relations over $O_F$. Put $q=hM$, $b_F=M\beta_h$; their triangular
basis there is $(t_F^{q-b_F},0)$ and $(-d,\beta)$ with
$v_F(\beta)=b_F$, for a uniformizer $t_F$. The semilinear action fixes
$v$ and sends $w$ to $w+v$. Applying $\sigma$ to the second relation
and subtracting $(\sigma\beta/\beta)$ times that relation gives
\[
 (\sigma-1)(d/\beta)-1\ \in\ t_F^{q-2b_F}O_F.
\]
Its trace is $-M$. If $d_F$ is the different exponent, trace duality
gives $\mathrm{Tr}(t_F^jO_F)=p^{\lfloor(j+d_F)/M\rfloor}\mathbb Z_p$:
this follows by testing inclusion of $p^{-i}t_F^jO_F$ in the inverse
different. Differentiating the cyclotomic polynomial and removing
the tame top extension of degree $p-1$ gives
$d_F=(h+1)M-1-(M-1)/(p-1)$. Hence
\[
 \beta_h\ge\frac h2-\frac{1-p^{-h}}{2(p-1)}.
\]
At full target precision $p^h$ the possible values of $v$ form the
ideal of valuation $\beta_h$. Indeed $W_h$ is killed by $p^h$,
its cyclic submodule generated by $v$ has normalized length
$h-\beta_h$, and linear forms into $O_h/p^h$ on that submodule
extend to $W_h$ by elementary-divisor duality over this principal
Artin ring. Thus full group precision and fixed precision have
different behavior even for one row.

There is a sharp limit on using several outputs. The second row
$r'_h=2e_h-r_h$ has the same primitive character valuations as $r_h$,
the same augmentation and compatibility. The proof just given applies
to it separately. But the lattice generated by $1,r_h,r'_h$ contains
$e_h$, since $2$ is a unit. As $x=x(1-e_h)$, its distinguished class
$v$ vanishes. Every normalized first-order map on the joined lattice
has zero speed. A local first-order map for a single output is
insufficient to construct a simultaneous arithmetic lift.
\section{Recovering an omitted character value by bounded Gauss averaging}
\label{sec:arch-gauss-averaging}
Let $\mathcal O$ be the ring of integers of a finite extension of
$\mathbb Q_p$ and take $v_p(p)=1$. Say that a power series is bounded
if its coefficients belong to $p^{-c}\mathcal O$ for one fixed $c$.
For a $p$-power root of unity $\zeta$ the associated character on
$\mathbb Z_p$ is $x\mapsto\zeta^x$.

\paragraph{A boundary interpolation criterion.}
Let $F(U,Z_1,\ldots,Z_d)$ and $G(Z_1,\ldots,Z_d)$ be bounded.
For each $m\ge1$ consider tuples $z=(\zeta_1-1,\ldots,\zeta_d-1)$
at which the successive orders $p^{n_i}$ are sufficiently large:
$n_1>N_1(m)$, $n_2>N_2(m,n_1)$, and so on. Assume the following
holds for \emph{all} such tuples and all nontrivial $\xi$ with
$\xi^{p^m}=1$:
\[
 F(\xi-1,z)=\sum_{j=1}^h g_j(z)\xi^{c_j(z)},\qquad
 G(z)=\sum_{j=1}^h g_j(z).
\]
Here $c_j(z)\in\mathbb Z_p$ may depend on $m,z$, but is constant
as $\xi$ varies. The $g_j(z)$ belong to the valuation ring of the
specialization after multiplication by a common $p^c$ independent
of $m,z$. Then $F(0,Z)=G(Z)$.

Indeed, the sum of $(\xi-1)^i$ over all $\xi^{p^m}=1$ is in
$p^m\mathbb Z_p$, by binomial expansion and character orthogonality.
The same estimate holds with loss at most $c$ for the sum of a
bounded series, first for polynomials and then by convergence at the
finite set of roots. The sums of the individual characters on the
right are either zero or $p^m$. Subtracting the two full sums leaves
only $F(0,z)-G(z)$; its valuation is at least $m-C$, with $C$ fixed.
For a nonzero bounded one-variable series $h$ over a finite
extension, Weierstrass preparation gives
\[
 \lim_{n\to\infty} v_p(h(\zeta_{p^n}-1))
       =\inf_i v_p([T^i]h).
\]
In fact, for sufficiently large $n$ the distinguished polynomial
of degree $e$ contributes $e/\varphi(p^n)$ to the valuation, since
its other coefficients have positive valuation; the unit contributes
zero. Apply this assertion successively to the last variable with
earlier ones fixed, then to each coefficient as a function of the
preceding variable. At a fixed $m$ all coefficients of
$F(0,Z)-G(Z)$ have valuation at least $m-C$. Taking arbitrarily
large $m$ gives the claim.

\paragraph{Stable Gauss sums.}
Let $H/\mathbb Q_p$ be unramified, write $O_H$ for its integers,
and fix a conductor-zero trace additive character $\psi$.
Use the convention
\[
 G(\lambda;p^N)=
 \sum_{a\in(O_H/p^N)^\times}\lambda(a)\psi(-a/p^N).
\]
Suppose $\lambda$ has conductor exponent $N\ge2s$, $s\ge1$.
The restriction of $\lambda$ to $1+p^{N-s}O_H$ has the form
$\lambda(1+p^{N-s}b)=\psi(d_\lambda b/p^s)$ for a unit
$d_\lambda$ modulo $p^s$, by the perfect additive trace pairing.
For every $\nu$ of conductor exponent at most $s$,
\[
 G(\lambda\nu^{-1};p^N)=\nu(d_\lambda)^{-1}G(\lambda;p^N).
\]
To prove this, sum first over the translates of
$1+p^{N-s}O_H$ modulo $p^N$. On the coset of $a$ the inner
additive sum is zero unless $a=d_\lambda\bmod p^s$.
All surviving cosets have the same $\nu(a)$. This gives a source
of the character sums used in the criterion above.

\paragraph{A ramified factor depending on the other characters.}
Over $\mathbb Q_p$ with $p\ge11$, choose a primitive root $g$ modulo
$p^6$, and write $[t]$ for the character with $[t](g)=
\exp(2\pi i t/p^5)$ and $[t](p)=1$. Put $\mu_1=[p^4]$, $\mu_2=[1]$,
$\chi=\mu_1\mu_2$. For $b=1,2$ take $\psi_1=[p^2]$,
$\tau_1=[bp^3]$, $\tau_2=[p]$, and
$\psi_2=\chi/\psi_1$, $\xi_i=\psi_i/\tau_i$. The conductor
exponents of $\mu_1,\tau_1,\psi_1,\tau_2,\mu_2$ are $2,3,4,5,6$.
At each occurrence write $G$ for the primitive Gauss sum at its
character's conductor; twisting a character of conductor at least
four by $\mu_1^{-1}$ preserves the conductor. Consider
\[
 R_b=\prod_{i=1}^2
 \frac{G(\mu_1^{-1}\xi_i)G(\mu_2^{-1}\xi_i)}
      {G(\xi_i)G(\chi^{-1}\xi_i)} .
\]
For a character $[p^{6-c}t]$ of conductor exponent $c\ge4$, the
parameter $d_\lambda$ modulo $p^2$ is
$t(\log_p(g)/p)^{-1}$. Here the logarithm is taken on the
principal-unit part of $g$. This follows by writing
$\log_p(1+p^{c-2}u)\equiv p^{c-2}u\bmod p^c$ and converting
to the exponent of $g$. For $\xi_1,\xi_2,\mu_2^{-1}\xi_1,
\mu_2^{-1}\xi_2$ the four normalized unit parameters $t$ modulo
$p^2$ are respectively $1-bp,1-p,-1,-(1+p)$.
The stable formula cancels the common logarithm factor and yields
\[
 R_b=\mu_1\left(\frac{1+p}{(1-bp)(1-p)}\right)
    =\mu_1(1+(b+2)p).
\]
The two values differ by a primitive $p$th root of unity.
In particular $R_b$ cannot be absorbed into a constant depending
only on $\mu_1,\mu_2$ and their conductor data. This tests this
particular four-character simplification on the generic ramified
stratum. For the interpolation setting see
\href{https://doi.org/10.1017/fms.2015.16}{Wan, Section 4.1.3,
Propositions 68 and 69}. Any application to a global family also
requires identification of the actual interpolated coefficient and
its lattices.

\paragraph{Provenance.} The three-corner calculation and its examples
come from entries 296, 304 and 318 of the BSD research record. The graph
calculation and source-pairing example come from the November 16--17
mixed-characteristic graph and source-Casimir analyses. The Gauss-row results are recovered from the first-level period calculation and its independent audit, the interlevel derivative comparison, and the higher-level fixed-precision argument. The final
precision theorem and its equality example are in the November 17
finite-precision checkpoint. Their intended application to a unitary
Hecke algebra still requires the arithmetic finite character, the full
matched coefficient class and the critical local conditions.
\end{document}
